% 
%  chapter6.tex
%  ISELthesis
%  
%  Created by Matilde Pós-de-Mina Pato on 2012/10/09.
%
\chapter{Case Study: Real-Time Fraud Detection for Card Payments}
\label{cha:case-study}

The previous chapters described the Function Registry and resolution cascade in isolation from
any particular domain (Chapter~\ref{ch:proposed_solution}) and traced their implementation
(Chapter~\ref{cha:impl}); Chapter~\ref{cha:evaluation} then measures what they cost. This
chapter grounds those mechanisms in a single, illustrative scenario drawn from retail banking:
a card-payment authorization workflow that screens every transaction for fraud before it is
settled. The scenario is not a validated production deployment; it is a realistic composition of
the \texttt{FraudPipeline} example already introduced in Section~\ref{sec:developer-workflow},
extended to make an authorization decision, used here
to walk through the integration end to end and to discuss where it fits an actual banking
environment, and where it does not yet fit one.

The chapter is organised as follows. Section~\ref{sec:case-scenario} sets out the business
requirements a payment-authorization workflow imposes and why they line up with the problem this
work addresses. Section~\ref{sec:case-design} defines the workflow in the OmniFlow DSL.
Section~\ref{sec:case-deployment} walks the resolution cascade through three deployments of that
same workflow (to a test and a production AWS account, then to GCP) to show how the bindings
follow the deployment target.
Section~\ref{sec:case-runtime} follows a transaction through the deployed workflow at runtime, and
Section~\ref{sec:case-discussion} discusses the scenario's relevance and its limits, connecting
them to the gaps assessed in Chapter~\ref{cha:conclusions}.

\section{Scenario and Requirements}
\label{sec:case-scenario}

Consider a bank that authorizes card payments through a workflow rather than a monolithic
service: each transaction is screened by two independent scoring functions before the bank
decides whether to settle it. \texttt{fraud-check} is a rules/ML-based service that flags
transactions matching known fraud patterns; \texttt{risk-score} computes a numeric risk score from
transaction and account features. Both are internal functions owned by the bank's fraud team,
deployed and iterated on independently of the orchestration logic. A third, external call settles
approved transactions against the bank's core banking system, which is a distinct system with its
own release cycle and is not managed through the registry.

Three requirements are characteristic of this setting and motivate the integration described in
this dissertation:

\begin{itemize}
    \item \textbf{One workflow, several accounts.} Suppose the bank keeps development, testing and
    production in separate cloud accounts (or projects), a common way to isolate environments, and
    promotes the workflow from one to the next. A function deployed in each account has a different
    endpoint there: its ARN contains the account ID. If the workflow definition embedded those
    endpoints, every promotion would require editing them by hand, and the functions would have to
    be deployed in the new account before the workflow could even be completed: the two problems
    that motivate the Function Registry (Section~\ref{sec:int_cha_int_uni_multi}).
    \item \textbf{An exit strategy.} The Digital Operational Resilience Act (DORA) requires EU
    financial entities to put in place exit strategies for ICT services supporting critical or
    important functions, and to assess ICT concentration risk before contracting such
    services~\cite[Arts.~28(8) and~29]{dora2022}. A fraud-screening workflow is a plausible
    candidate for such a function. Exiting a provider is easier if the orchestration and the
    functions it calls can be deployed elsewhere, not just the infrastructure underneath them.
    \item \textbf{Auditability of bindings.} A bank's change-management process typically requires
    a record of which endpoint each deployment was bound to, which the
    registry's \texttt{updatedAt}-stamped entries provide in embryonic form
    (Section~\ref{subsec:registry-evolution}), even though a production deployment would need the
    stronger guarantees discussed in Section~\ref{sec:case-discussion}.
\end{itemize}

\section{Workflow Design}
\label{sec:case-design}

Listing~\ref{lst:case-payment-workflow} defines \texttt{PaymentAuthorization}. The first two
steps call \texttt{fraud-check} and \texttt{risk-score} through \texttt{internalFunction}, so
their endpoints are resolved through the registry at deployment time rather than hard-coded, as
described in Section~\ref{subsec:endpoint-resolution}. Unlike the partner-owned risk-scoring API
of Listing~\ref{lst:wf-external-part}, which the workflow reaches through a literal endpoint, both
scoring calls here run on functions the bank owns, and are therefore internal and
registry-managed. Their results are captured in \texttt{fraudResult} and \texttt{riskResult}. A third step,
\texttt{decision}, is a conditional (\texttt{switch}) step: it jumps to
\texttt{decline-transaction} if the fraud check flagged the transaction or the risk score exceeds
a threshold, and otherwise falls through to the \texttt{default}, \texttt{approve-transaction}.
The approve path issues an external call to the bank's own settlement endpoint, declared with a
literal \texttt{host}/\texttt{path}, exactly like \texttt{external-risk-api} in
Listing~\ref{lst:wf-external-part}.

\begin{lstlisting}[language=Kotlin,caption={PaymentAuthorization workflow: two internal scoring calls feeding a decision.},label={lst:case-payment-workflow},basicstyle=\ttfamily\small,breaklines=false]
val PaymentAuthorization = workflow {
  name("PaymentAuthorization")
  description(
    "Real-time fraud screening for card payment authorization")
  input("transaction")
  steps(
    step {
      name("fraud-check")
      context(call {
        method(POST)
        internalFunction("fraud-check", "./deploy/fraud-check.json")
        body(variable("transaction"))
        result("fraudResult")
      })
    },
    step {
      name("risk-score")
      context(call {
        method(POST)
        internalFunction("risk-score", "./deploy/risk-score.json")
        body(variable("transaction"))
        result("riskResult")
      })
    },
    step {
      name("decision")
      context(switch {
        conditions(
          condition {
            match(variable("fraudResult.body.flagged")
              equalTo value(true))
            jump("decline-transaction")
          },
          condition {
            match(variable("riskResult.body.score")
              greaterThan value(75))
            jump("decline-transaction")
          }
        )
        default("approve-transaction")
      })
    },
    step {
      name("approve-transaction")
      context(call {
        method(POST)
        host("core-banking.internal.example")
        path("/v1/settlements")
        body(variable("transaction"))
        result("settlementResult")
      })
    },
    step {
      name("decline-transaction")
      context(assign {
        variables(variable("decision") equalTo value("declined"))
      })
    }
  )
  result("decision")
}
\end{lstlisting}

Nothing in Listing~\ref{lst:case-payment-workflow} names a cloud provider, a region or a concrete
endpoint for either scoring function: that information is supplied once, at deployment time,
through the QuickFaaS descriptors referenced by \texttt{internalFunction} and the choice of
provider deployer. This is what makes the same definition deployable to either provider, as
Section~\ref{sec:case-deployment} shows next.

\section{Deploying Across Accounts and Providers}
\label{sec:case-deployment}

Suppose neither \texttt{fraud-check} nor \texttt{risk-score} has been deployed yet, and the bank
deploys \texttt{PaymentAuthorization} to its AWS test account first. The cascade runs as in the
walkthrough of \S\ref{sec:walkthrough}, and reaches Level~3 for both calls: QuickFaaS deploys \texttt{fraud-check} and \texttt{risk-score}
as Lambda functions from the descriptors referenced in the workflow, following the AWS deployment
sequence of Figure~\ref{fig:aws-deploy-sequence}. Both ARNs are
written to the registry, and \texttt{PaymentAuthorization} is deployed as a state machine.

The bank then promotes \texttt{PaymentAuthorization} to its production account. The workflow
definition does not change; only the descriptors do, since they name the production account's
deployment bucket (and, optionally, its execution role). The production deployment uses its own
registry file (the deployers accept a registry path, so one per account is a configuration
choice): OmniFlow bootstraps it from the production account, finds neither scoring function, and
Level~3 deploys both there, recording ARNs that carry the production account ID. No endpoint is
copied between accounts. The registry file must not be shared between the two accounts: an entry
written for the test account names a function the production account does not have, so its
validation fails, the entry is treated as stale, and the production deployment aborts
(\S\ref{sec:walkthrough}).

Now suppose the bank's exit strategy calls for the same fraud-screening workflow to run on a
second provider (Section~\ref{sec:case-scenario}). Deploying it to GCP takes GCP descriptors (targeting Cloud Functions instead of Lambda) and the
GCP deployer. On this first GCP deployment, the GCP registry is
bootstrapped separately from the AWS one, both scoring functions miss at Level~1 and Level~2, and
Level~3 deploys them as first-generation Cloud Functions instead, following the ZIP-based
deployment path described in \S\ref{sec:background_quickfaas}. The two registries stay
independent, one per provider (\texttt{function-registry.aws.json} and
\texttt{function-registry.gcp.json} by default), so the same \texttt{functionRef} resolves to a
Lambda ARN under one deployment and to a Cloud Functions URL under the other. On GCP, however, those
bindings are not validated on later deployments (\S\ref{subsec:store-internals}). The uniform
\texttt{url} field discussed in
Section~\ref{subsec:registry-evolution} is precisely what makes this substitution transparent to
the workflow definition.

\section{Runtime Invocation}
\label{sec:case-runtime}

Once deployed, the two deployments of \texttt{PaymentAuthorization} differ only in how the two
scoring calls are invoked. On AWS, Step Functions invokes \texttt{fraud-check} and
\texttt{risk-score} through the native \texttt{lambda:invoke} integration of
Section~\ref{sec:runtime-invocation}; on GCP, they are HTTP calls to the resolved Cloud Functions
URLs. The rest of the workflow is the same on both: the \texttt{decision} step becomes a native
conditional (a \texttt{Choice} state on AWS), and \texttt{approve-transaction} an ordinary
external call.

\section{Discussion}
\label{sec:case-discussion}

This scenario exercises the integration's central claim in a setting where it plausibly matters:
a bank can promote the same workflow definition across its accounts and, as part of an exit
strategy, deploy it on a second cloud provider, changing deployment descriptors rather than the
workflow itself. In each new target the missing functions are deployed and their endpoints bound
without anyone copying an ARN or URL by hand. The fraud team's frequent model updates are a
different matter: a retrained \texttt{fraud-check} redeployed in place keeps its endpoint
(\S\ref{sec:resolution-cascade}), so the deployed workflow calls the new model unchanged. Conversely, a workflow deployment never deploys the new model itself, because the
cascade only deploys missing functions (Section~\ref{sec:concl-critical}).

The scenario equally exposes why this remains illustrative rather than production-ready: the
unguarded registry and the broad, automatically granted IAM role assessed in
Section~\ref{sec:concl-critical} are not controls a bank's change-management or security function
would accept for a function that screens payment data.
The hardening required before the workflow in Listing~\ref{lst:case-payment-workflow} could
authorize real transactions is set out in Section~\ref{sec:concl-future}.






















































































